\section{Deterministic Threshold Signature Schemes}

We begin with the definition of a deterministic threshold signature scheme,
and then define the notion of unforgeability.
We build upon standard definitions and notions of unforgeability for threshold signatures in the literature~\cite{CritesKM23,GennaroJKR01,GennaroRJK07,GennaroG19}.

Intuitively, our definition for deterministic threshold signature schemes is
identical to that of randomized threshold signature schemes,
with the exception that in the deterministic setting,
the signing algorithms are deterministic and the signer does not maintain state between rounds.
Furthermore,
each signing round in the deterministic setting is given as input the message $m$,
coalition of signers $\coalition$,
and secret signing key share $\sk_i$.

\begin{definition}[Deterministic Threshold Signatures]
  A \defn{deterministic threshold signature scheme} $\dthreshsig$ with an interactive signing protocol consisting of $\numrounds$ rounds is a tuple of PPT algorithms
$\threshsig = (\setup, \keygen, \allowbreak (\sign_1, \ldots, \sign_\numrounds), \allowbreak \combine, \allowbreak \verify)$, defined as follows:


\begin{itemize}

  \item $\setup(1^\lambda) \rightarrow \pp$:
A probabilistic algorithm that accepts as input a security parameter $\lambda$ and outputs public parameters $\pp$,
which are then implicitly given as input to all other algorithms.

  \item $\keygen(n,\thresh,\minsigners) \rightarrow (\pkgroup, \{ \pk_i, \sk_i\}_{i \in \set{n}})$:
A probabilistic algorithm that takes as input the total number of possible signers $n$,
the corruption threshold $\thresh$,
and the minimum number of participating signers $\minsigners$.
Outputs the public key $\pkgroup$ representing the set of $n$ signers,
the set $\{ \pk_i, \sk_i \}_{i \in \set{n}}$ of public and secret key shares for each signer.
%It is assumed that participant $i$ is sent its respective share $\sk_i$ privately.
%\ian{Nowhere in the paper do we use the $\pk_i$. Do we want to keep them?}
%\chelsea{it is fairly standard at this point to include them...}

  \item $(\sign_1, \ldots, \sign_\numrounds) \rightarrow \{\tsigmsg_1, \ldots \tsigmsg_\numrounds\}_{k \in \coalition}$:
A set of deterministic algorithms where each algorithm represents a single stage in an interactive signing protocol
performed by signing party $\partyid \in \set{n}$ in a signing set $\coalition \subseteq \set{n},
|\coalition| \geq \minsigners$ with respect to a message $\msg$, defined as follows:

% Ian compressed this to remove unnecessary whitespace; the original is
% below.
    \(\tsigmsg_1 \gets \sign_1(\partyid, \sk_\partyid, \msg, \coalition),
    \quad \ldots, \quad
    \tsigmsg_\numrounds \gets \sign_\numrounds(\partyid, \sk_\partyid,
    \msg, \coalition, \{\tsigmsgin_{\numrounds-1}\}_{i \in
    \coalition})\)

%  \begin{align*}
%    \tsigmsg_1 &\gets \sign_1(\partyid, \sk_\partyid, \msg, \coalition)  \\
%    %\tsigmsg_2 &\gets \sign_2(\partyid, \sk_\partyid, \msg, \coalition, \{\tsigmsgin_1\}_{i \in \coalition}) \\
%    \vdots \\
%    \tsigmsg_\numrounds &\gets \sign_\numrounds(\partyid, \sk_\partyid, \msg, \coalition, \{\tsigmsgin_{\numrounds-1}\}_{i \in \coalition})
%  \end{align*}
where $\tsigmsg_1, \ldots, \tsigmsg_\numrounds$ are protocol messages produced by party $\partyid \in \coalition$.

  \item $\combine(\pkgroup, \msg, \coalition, \{(\tsigmsgin_1, \ldots, \tsigmsgin_\numrounds) \}_{i \in \coalition}) \rightarrow \sigma$:
  A deterministic algorithm that takes as input the public key $\pk$, the message
    $\msg$, the set of signers $\coalition$, and the set of protocol messages
  sent by each party during the $\sign_1, \ldots, \sign_\numrounds$ signing stages, and outputs a joint signature $\sigma$.

  \item $\verify(\pkgroup, \msg, \sigma) \rightarrow 0/1$:
A deterministic algorithm that takes as input the public key $\pkgroup$,
a message $\msg$,
and signature $\sigma$ and outputs $1$ to indicate accept if the signature verifies;
otherwise, it outputs $0$ to indicate reject.
\end{itemize}
\end{definition}

\begin{remark}[Distributed key generation]\label{rmk:dkg}
  Our definition assumes a centralized key generation algorithm $\KeyGen$ to
  generate the public key $\pkgroup$ and public key shares $\{ \pk_i, \sk_i \}_{i \in \set{n}}$.
  However, our scheme and proofs can be adapted to use a fully decentralized distributed key generation protocol (DKG).
\end{remark}

\begin{remark}[Deferring the Choice of Coalition to Later Rounds]
  Our definition assumes that the coalition of signers $\coalition$ is provided in the first round of signing $\Sign_1$.
  However, some constructions (as is the case with \glitter) may defer the choice of $\coalition$ to later rounds.
\end{remark}

\paragraph{Correctness.} A deterministic threshold signature scheme $\dthreshsig$ is \emph{correct} if for all security parameters $\lambda$,
all $1 \leq \thresh \leq \minsigners \leq n$,
all participant identifiers $k \in \set{n}$,
all $\coalition \subseteq \set{n}$ such that $\minsigners \leq | \coalition | \leq n$,
and for all messages $m \in \{0,1\}^*$,
the following relation holds:

  \begin{align*}
    \dthreshsig.\verify(\pkgroup, \msg, \sigma) = 1, \text{ for } \\
    (\pkgroup, \{ \pk_i, \sk_i \}_{i \in \set{n}}) \rgets \dthreshsig.\keygen(n, \thresh,\minsigners),
    \text{ where}\\
    \tsigmsg_{1} \gets \dthreshsig.\sign_1(\partyid, \sk_\partyid, \msg, \coalition),
    \ldots,
    \tsigmsg_{\numrounds } \gets \dthreshsig.\sign_\numrounds(\partyid, \sk_\partyid, \msg, \coalition, \{\tsigmsgin_{\numrounds-1} \}_{i \in \coalition}), \text{ and} \\
    \sigma \gets \dthreshsig.\combine(\pkgroup, \msg, \coalition, \{ \tsigmsgin_{1}, \ldots, \tsigmsgin_{\numrounds } \}_{i \in \coalition})
  \end{align*}

\begin{figure}[t]
  \centering
	%\begin{linedgamespec}
	\begin{pchstack}[boxed]
		\begin{pcvstack}
		\procedure{ $\UnforgExp_{\dthreshsig,\adv}(\lambda, n, t,\minsigners)$ }{
			\pp \rgets \dthreshsig.\setup(1^\lambda) \\
      Q \gets \emptyset \pccomment{set of queried messages} \\
      (\corrupt, \advstate) \rgets \adv(\pp, n, t, \minsigners) \\
      \pcif \lvert \corrupt \rvert \geq t \\
      \pcind \pcreturn \bot \\
      \honest \gets \set{n} \setminus \corrupt \\
      (\pkgroup, \{ \pk_i, \sk_i \}_{i=1}^n) \rgets \dthreshsig.\keygen(n,\thresh,\minsigners) \\
      \mathsf{in} \gets (\pk, \{ \pk_i \}_{i =1}^n,  \{\sk_j \}_{j \in \corrupt},  \advstate)  \\
      (\msg^*, \sigma^*) \rgets \adv^{\BigO^{\tiny{\Sign_i, i \in \set{\numrounds}}}}(\mathsf{in}) \\
      \pcif \msg^* \notin Q \wedge \dthreshsig.\verify(\pkgroup, \msg^*, \sigma^*) = 1\\
      \pcind \pcreturn 1 \\
			\pcreturn 0
		}
   		\pcvspace
		\end{pcvstack}
    		\pchspace
		\begin{pcvstack}
			\procedure{ $\BigO^{\Sign_1}(\partyid, \msg, \coalition)$ }{
        \pccomment{$\partyid$ denotes the participant identifier} \\
				\tsigmsg_1 \gets \dthreshsig.\Sign_1(\partyid, \sk_\partyid, \msg, \coalition) \\
        \pcreturn \tsigmsg_1
			}
        \vdots
      \procedure{ $\BigO^{\Sign_j}(\partyid, \msg, \coalition, \{\tsigmsgin_{j-1} \}_{i \in \coalition})$ }{
        \text{ \algcom for $j \in \{ 2, \ldots, \numrounds-1\}$} \\
				\tsigmsg_j \gets \dthreshsig.\Sign_i(\partyid, \sk_\partyid, \msg, \coalition, \{\tsigmsgin_{j-1} \}_{i \in \coalition}) \\
        \pcreturn \tsigmsg_j
			}
        \vdots
        \procedure{$\BigO^{\Sign_\numrounds}(\partyid, \msg, \coalition, \{\tsigmsgin_{\numrounds-1} \}_{i \in \coalition})$}{
				Q \gets Q \cup \{\msg\} \\
        \tsigmsg_\numrounds \gets \dthreshsig.\Sign_\numrounds(\partyid, \sk_\partyid, \msg, \coalition, \{ \tsigmsgin_{\numrounds-1} \}_{i \in \coalition}) \\
        \pcreturn \tsigmsg_\numrounds
		}
		\end{pcvstack}
	\end{pchstack}
	%\end{linedgamespec}
\caption{Unforgeability experiment for a deterministic threshold signature scheme with a $\numrounds$-round signing protocol.
  }
\label{fig:euf-cma-threshold-deterministic}
\end{figure}
%






\subsection{Unforgeability}
%
We present a game-based definition of unforgeability for a deterministic threshold signature scheme in
Figure~\ref{fig:euf-cma-threshold-deterministic}.
This notion of unforgeability is analogous to the standard notion of chosen message attack (EUF-CMA) for standard signature schemes~\cite{GoldwasserMR88}.
We present the adversary as a static adversary,
which is allowed to corrupt at most $\thresh-1$ signers.

The game assumes a static adversary that picks the players it corrupts at the beginning of the game.
The public parameters $\pp$ are implicitly given as input to all algorithms,
and $\tsigmsg_1, \ldots, \tsigmsg_\numrounds$ represent protocol messages
sent by participant $\partyid$ throughout the interactive signing protocol.
%
The static unforgeability game for a deterministic threshold signature scheme is nearly identical to that of a randomized scheme,
with the exception that the environment does not maintain state between invocations of signing oracles by the adversary,
and the adversary provides as input the message $\msg$ and coalition of signers $\coalition$ (of its choosing) for each call to a signing oracle.

In the static unforgeability game for a deterministic threshold signature scheme,
the environment allows the adversary to sample the corrupt parties $\corrupt \subset \set{n}$.
If the set of corrupt parties is smaller than the corruption threshold $t$,
it derives the set of honest parties $\honest \subseteq \set{n}$ as the remaining parties in $\set{n}$.
The environment then runs $\keygen$ to derive the joint public key $\pkgroup$ representing the set of $n$ signers,
the set of public key shares $\{ \pk_i \}_{i \in \set{n}}$,
and the secret key shares $ \{ \sk_i \}_{i \in \set{n}}$.
The adversary is then run on input $n, t, \minsigners, \pp, \pkgroup$, $\{ \pk_i \}_{i \in \set{n}}$,
and the corrupt signing key shares $ \{ \sk_j \}_{j \in \corrupt}$.
%
The adversary can then query any honest signers $k \in \honest$ of its choosing at each step in the signing protocol,
%by querying oracles $\BigO^{\Sign_1}, \ldots, \BigO^{\Sign_\numrounds}$,
and has full power over choosing the set of signers $\coalition$ and the message $\msg$.
Additionally, the adversary may query the signing round oracles in arbitrary order.
%in particular that unlike for randomized stateful signature schemes,
%This gives substantially more power to the adversary than in the randomized stateful case.
%\chelsea{this is also possible in the randomized steting.}
However,
unlike in the randomized setting as shown in Figure~\ref{fig:euf-cma-threshold},
the environment for a deterministic threshold signature scheme does not maintain any session identifiers,
or state about ongoing signing sessions.
%
The adversary wins if it can produce a valid forgery $\sigma^*$ with
respect to the joint public key $\pkgroup$ on a message $\msg^*$ that has not been queried to $\oracle^{\sign_\numrounds}$ (i.e., in the final round of signing).
%
Importantly, this definition allows the adversary to be \emph{rushing},
meaning it can wait to produce its outputs after having seen the honest outputs first.
The adversary may also be \emph{concurrent},
meaning it can open simultaneous signing sessions at once,
or choose not to complete a signing session.

\begin{definition}[Unforgeability]
Let the advantage of a static adversary $\adv$ playing the unforgeability game
  against a deterministic threshold signature scheme $\dthreshsig$ as defined in Figure~\ref{fig:euf-cma-threshold-deterministic} be as follows:

\[
  \advant{uf}{\dthreshsig, \adv}(\lambda, n, \thresh,\minsigners) = \Pr[ \UnforgExp_{\dthreshsig,\adv}(\lambda,n,t,\minsigners)= 1]
\]
%
A deterministic threshold signature scheme $\dthreshsig$ is \emph{unforgeable} if for all PPT adversaries $\adv$,
$\advant{uf}{\dthreshsig,\adv}(\lambda, n, t,\minsigners) $ is negligible
in $\lambda$,
for all $n, t,\minsigners \in \mathbb{N}$,
where $t \leq \minsigners \leq n$.
\end{definition}
